\lecture{12}{Thu 14 Nov 2024 09:31}{More Rings}
\begin{remark}
	The $p$-adic fields have a strange topology where if two open sets intersect then they are equal.
\end{remark}

\begin{definition}[Unity]\label{dfn:UN}
	A ring $R$ has unity if there is $1\in R$ such that $a1=1a=a$ for all $a\in R$.
\end{definition}
\begin{note}
Some books only talk about rings with unity, and call a ring without unity a rng. A lot can be said about rings with unity that do not generalize. However, we are not doing so.
\end{note}
\begin{definition}[Integral Domain]\label{dfn:ID}
	A commutative ring $R$ with unity is an integral domain if for every $a,b\in R$ such that $ab=0$, then either $a=0$ or $b=0$. We say there are no zero divisors.
\end{definition}
\begin{definition}[Domain]\label{dfn:D}
	A commutative ring is a domain if there are no zero divisors.
\end{definition}
\begin{definition}[Field]\label{dfn:Field}
	A field is a commutative ring with unity, no zero divisors, and for all nonzero $a\in R$ there exists $a^{-1}\in R$ such that $aa^{-1}=1$.
\end{definition}
\begin{definition}[Division Ring]\label{dfn:Div}
	A division ring is a ring with unity and no nonzero divisors.
\end{definition}
We are only concerned about the following categories.
\begin{center}
	\begin{tikzpicture}
		\node (A) at (0,0) {Ring};
		\node (B) at (3,-1.5) {Ring with Unity};
		\node (C) at (-3,-1.5) {Commutative Ring};
		\node (D) at (0,-3) {Integral Domain};
		\node (E) at (0,-4.5) {Field};
		\draw (A) -- (B) -- (D) -- (E);
		\draw (A) -- (C) -- (D);
	\end{tikzpicture}
\end{center}

\begin{definition}[Unit]\label{dfn:Unit}
	Let $R$ be a ring with $a\in R\setminus\left\{ 0 \right\} $. If $a$ has a multiplicative inverse, we say $a$ is a unit of $R$. We denote the set of units in $R$ as $U(R)$.
\end{definition}
\begin{definition}[Zero Divisor]\label{dfn:ZD}
	A nonzero element $a\in R$ in a commutative ring is called a zero divisor if there is a nonzero element $b\in R$ such that $ab=0$.
\end{definition}
\begin{remark}
	Proof by googling it: a noncommutative ring with unity can also have zero divisors. A rng has left/right zero divisors.
\end{remark}
Recall that $M_2(\mathbb{C})$ is a ring. It only has unity though. We want to find a matrix ring that has some properties. Let
\[
	1=I_2,\qquad \hati=\begin{pmatrix}
		0&1\\
		-1&0
	\end{pmatrix},\qquad \hatj=\begin{pmatrix}
		0&i\\
		i&0
	\end{pmatrix} ,\qquad \hatk=\begin{pmatrix}
		i&0\\
		0&-i
	\end{pmatrix}  
.\]
Define
\[
	\mathbb{H}=\left\{ a1+b\hati+c\hatj+d\hatk \mid a,b,c,d\in\mathbb{R} \right\} 
.\]
We call $\mathbb{H}$ the set of \textbf{quaternions}. We will show $\mathbb{H}$ is a non-commutative ring with unity, but without zero divisors. Recall that
\[
	\hati^2=\hatj^2=\hatk=\hati\hatj\hatk=-1,
\]
\[
	\hati\hatj=\hatk,\quad \hatj\hatk=\hati,\quad\hatk\hati=\hatj
,\]
\[
	\hatj\hati=-\hatk,\quad\hatk\hatj=-\hati,\quad\hati\hatk=-\hatj
.\]
We say $\mathbb{H}$ is anti-commutative. Note that if we multiply a quaternion $\mathbf{h}$ with its conjugate $\overline{\mathbf{h}} $, we have
\[
	\left( a1+b\hati+c\hatj+d\hatk \right) \left( a1-b\hati-c\hatj-d\hatk \right) =a^2+b^2+c^2+d^2
.\]
Therefore, these two things multiply to zero if and only if $a,b,c,d=0$. We have now shown for $\mathbf{h}\neq 0$, it and it's conjugate can be divided by this value. If $a1+b\hati+c\hatj+d\hatk\neq 0$, then
\[
	\left( a1+b\hati+c\hatj+d\hatk \right) \left( \frac{a1-b\hati-c\hatj-d\hatk}{a^2+b^2+c^2+d^2} \right) =1
.\]
This tells us that every element that is nonzero is a unit. It remains to be shown \textbf{that units are not zero divisors.}
\begin{theorem}[Properties of Rings]\label{thm:PRG}
	Let $R$ be a ring, and let $a,b\in R$. Then
	\begin{itemize}
		\item $a0=0a=0$.
		\item $a(-b)=(-a)b=-(ab)$.
		\item $(-a)(-b)=ab$.
		\item $a(b-c)=ab-ac$.
	\end{itemize}
	If $R$ has unity, then
	\begin{itemize}
		\item $(-1)a=-a$.
		\item $(-1)(-1)=1$.
		\item The unity 1 is unique.
	\end{itemize}
\end{theorem}
\begin{proof}
	(1)
	\[
		a0=a(0+0)=a0+a0
	.\]
	Since there are inverses under addition, $a0-a0=a0+a0-a0$ implies $a0=0$. Same for $0a$.

	(2)
	\[
		a(-b)+ab=a\left( -b+b \right) =a0=0
	.\]
	Since inverses are unique under addition, $a(-b)=-ab$. Similarly,
	\[
		(-a)b+ab=(-a+a)b=0b=0
	.\]
	Therefore, $a(-b)=(-a)b=-ab$.

	(3)
	\[
		(-a)(-b)+((-a)b)=(-a)(-b+b)=(-a)0=0
	.\]
	Since inverses are unique under addition, $(-a)(-b)=-((-a)b)=-(-ab)=ab$.

	(4)
	\[
		a(b-c)=a(b+(-c))=ab+a(-c)
	.\]
	By property (2), $a(-c)=-ac$. Therefore, $a(b-c)=ab-ac$.

	(5)
	\[
		a+(-1)a=1a-1a=(1-1)a=0a=0
	.\]
	Since inverses are unique under addition, $(-1)a=-a$.

	(6)
	\[
		(-1)(-1)-1=(-1)(-1)+(-1)1=(-1)(-1+1)=(-1)0=0
	.\]
	Since inverses are unique under addition, $(-1)(-1)=-(-1)=1$.

	(7) Let $e_1,e_2$ be identities in $R$. Since $e_1$ is an identity, $e_2=e_1e_2$. Since $e_2$ is an identity, $e_1=e_1e_2$. Therefore $e_1=e_2$.
\end{proof}
\begin{proposition}
	Let $R$ be a ring, and let $a\in R$ be a unit. Then $a$ is not a zero divisor.
\end{proposition}
\begin{proof}
	Let $a\in R$ be a unit, and let $ab=0$. Then $a^{-1}ab=a^{-1}0$. By properties of rings, $a^{-1}0=0$. Furthermore, $a^{-1}ab=1b=b$. Therefore, $b=0$, and $a$ is not a zero divisor.
\end{proof}
\begin{definition}[Subring]\label{dfn:Subring}
	Let $(R,+,\cdot )$ be a ring with subset $S\subseteq R$. Then $S$ is called a subring if $(S,+,\cdot )$ is a ring under the inherited operations.
\end{definition}
\begin{proposition}
	Let $R$ be ring with $S\subseteq R$. Then $S$ is a subring if and only if
	\begin{enumerate}
		\item $S\neq \varnothing$,
		\item $rs\in S$ for all $r,s\in S$,
		\item $r-s\in S$ for all $r,s\in S$.
	\end{enumerate}
\end{proposition}
\begin{proof}
	h
\end{proof}
For example, consider the ring $M_2(\mathbb{R})$, and condiser the subset $S\subseteq M_2(\mathbb{R})$ defined as
\[
	S\coloneqq \left\{ \begin{pmatrix}
			a&b\\
			0&c
	\end{pmatrix}  \mid a,b,c\in\mathbb{R} \right\} 
.\]
We will show $S$ is a subring of $M_2(\mathbb{R})$. Note that $I\in S$, so $S\neq \varnothing$. Furthermore, note that
\[
	\begin{pmatrix}
		a&b\\
		0&c
	\end{pmatrix} \begin{pmatrix}
		x&y\\
		0&z
	\end{pmatrix} =\begin{pmatrix}
		ax&ay+bz\\
		0&cz
	\end{pmatrix} \in S
.\]
Finally, note that
\[
	\begin{pmatrix}
		a&b\\
		0&c
	\end{pmatrix}- \begin{pmatrix}
		x&y\\
		0&z
	\end{pmatrix} =\begin{pmatrix}
		a-x&b-y\\
		0&c-z
	\end{pmatrix} 
.\]
Therefore, $S$ is a subring of $M_2(\mathbb{R})$.
\begin{remark}
	Let $R$ be a ring with unity. Note that $\left\{ 0 \right\} \subseteq R$ is a ring without unity. Therefore, subrings do not need to have the unity.
\end{remark}
Another good example of a ring is the Gaussian Integers.
\begin{definition}[Gaussian Integers]\label{dfn:GI}
	The gaussian integers are the set $\mathbb{Z}[i]$ defined as
	\[
		\mathbb{Z}[i]\coloneqq \left\{ a+bi \mid a,b\in\mathbb{Z} \right\} \subset \mathbb{C}
	.\]
\end{definition}
We claim $\mathbb{Z}[i]$ is an integral domain, but not a field. Let $\alpha =a+bi\in\mathbb{Z}\left[ i \right] $, and suppose $\alpha $ is a unit. Then there exists $\beta =c+di\in\mathbb{Z}[i]$ such that $\alpha \beta =1$. We want to show this does not happen in general. Consider the conjugate $\overline{\alpha} \overline{\beta } =1$. We then have $\alpha \beta \overline{\alpha } \overline{\beta } =1$. We have
\[
	\left( a+bi \right) \left( a-bi \right) \left( c+di \right) \left( c-di \right) =\left( a^2+b^2 \right) \left( c^2+d^2 \right) =1
.\]
We now have the product of two positive integers are $1$. This happens if and only if both integers are $1$. That is, $a^2+b^2=c^2+d^2=1$. Then $a=\pm1$ and $b=0$, or vice versa. These are the only cases where $\alpha $ has a unity $\beta $. Therefore, $\mathbb{Z}[i]$ is \textbf{not} a field.
\chapter{Integral Domains}
